\section{二次复杂度与推理时的显存账}
\label{sec:14-Deep-Learning/05-Attention-and-Sequence-Models/05}

你把上下文窗口从 4k 改到 32k，训练脚本立刻 OOM。把 batch 从 8 砍到 1，还是 OOM。这时候部署那边也来了消息：同一个模型在线上没有崩在启动，而是在生成到第三千个 token 左右的时候显存一路爬升然后崩掉，重启后又能跑一会儿。

同一个模型，同一段代码，训练和推理撞上的是两堵不同的墙。要把窗口做长，得先把两本账分开算。

\subsection{训练侧：\(n^2\) 那一项吃掉了一切}

\hyperref[sec:14-Deep-Learning/05-Attention-and-Sequence-Models/01]{``Self-Attention 是内容决定的加权读取''}已经指出分数矩阵是 \(n\times n\) 的；训练时的账要在此之上再乘三个因子——批量、层数、头数——因为反向传播需要把前向的注意力权重留着。

一层一个头的注意力权重矩阵占 \(n^2\) 个数。批量 \(B\)、层数 \(L\)、头数 \(h\) 的模型，光是这一项就要

\[
B\cdot L\cdot h\cdot n^2 \times (\text{每数字节数}).
\]

取 \(B=1\)、\(L=32\)、\(h=32\)、fp16。\(n=4096\) 时是 \(32\cdot32\cdot4096^2\cdot2\approx 34\) GB——已经放不下了，实践中靠重计算和分块把它压回去。\(n\) 翻到 32768，这一项直接乘 64。把 batch 砍成 1 完全不解决问题：\(B\) 只是个乘数，\(n^2\) 才是指数上的东西。

计算量同样是平方：\(QK^{\trans}\) 与 \(AV\) 各是一次 \(O(n^2 d)\) 的矩阵乘。而模型的其他部分——线性投影、前馈层——都是 \(O(n d^2)\)，随 \(n\) 线性。两者的比值是 \(n/d\)：\(n\) 远小于 \(d\) 时注意力只占很小一块，\(n\) 超过 \(d\) 之后注意力开始主导，之后每翻一倍长度就四倍开销。\(d=4096\) 的模型在 \(n=4096\) 附近正好是分水岭，这不是巧合——现有模型的窗口长度大多就卡在这个数量级上。

\subsection{推理侧：瓶颈换成了 KV cache}

自回归生成时上面那本账不成立了。每一步只有一个新 token 作为 Query，分数是一个 \(1\times n\) 的向量而不是 \(n\times n\) 的矩阵，注意力权重用完即弃。二次项消失了。

换上来的是缓存。历史位置的 Key 与 Value 在后续每一步都要用，重算它们等于把整个前缀再前向一遍，所以必须留在显存里。留多少？每个 token、每一层要存一组 \(K\) 和一组 \(V\)，宽度都是 \(d\)：

\[
\text{KV cache} \;=\; 2 \times L \times n \times d \times (\text{每数字节数}).
\]

代进一个七十亿参数量级的常见配置：\(L=32\)、\(d=4096\)、fp16。每个 token 每层要存 \(2\times 4096\times 2=16\) KB，三十二层就是 \(0.5\) MB。一个 token 半兆——这个数字值得记住，因为它把后面的估算都变成了心算。

生成 4096 个 token，缓存 2 GB；32768 个 token，缓存 16 GB。而这个模型的权重在 fp16 下是 14 GB。上下文一长，缓存比模型本身还大，而且它是随生成过程持续增长的，不像权重那样是常数。线上那个``跑一会儿才崩''的现象至此完全解释清楚了：显存不是被一次性占满的，是被一个 token 一个 token 喂满的。多路并发时每条请求各有一份缓存，占用还要再乘上并发数。

还有一件事容易被忽略。解码每一步要把整个 cache 从显存读一遍，只为算一个 \(1\times n\) 的分数向量和一次加权求和——读进来的数据量与做的浮点运算量几乎成正比，算术强度接近 1。这意味着解码阶段是访存受限的，GPU 的浮点峰值基本用不上，速度由显存带宽决定。这个观察在下面讨论 FlashAttention 时是关键。

\begin{center}
\begin{tikzpicture}[scale=0.95]
\draw[->,>=stealth] (0,0) -- (7.4,0);
\draw[->,>=stealth] (0,0) -- (0,4.5);
\node[anchor=north] at (3.2,-0.62) {\footnotesize 序列长度 \(n\)};
\node[rotate=90,anchor=south] at (-0.45,2.2) {\footnotesize 资源占用};
% 显存上限
\draw[dashed,black!60] (0,3.0) -- (7.0,3.0);
\node[anchor=south east] at (7.0,3.0) {\footnotesize 可用显存};
% n^2 曲线
\draw[line width=1pt,domain=0:6.1,samples=80] plot (\x,{0.08*\x*\x});
\node[anchor=south east] at (3.95,2.45) {\footnotesize 训练：注意力 \(O(n^2)\)};
% 线性曲线
\draw[line width=1pt,dash pattern=on 3pt off 2pt,domain=0:7.0,samples=2] plot (\x,{0.30*\x});
\node[anchor=north west] at (4.6,1.45) {\footnotesize 推理：KV cache \(O(n)\)};
% 交点
\fill (6.124,3.0) circle (1.6pt);
\draw[dotted,black!70] (6.124,0) -- (6.124,3.0);
\node[anchor=north] at (6.124,-0.05) {\footnotesize \(n_{\text{train}}\)};
\node[anchor=west,align=left] at (0.35,3.85) {\footnotesize 长度翻一倍，\\[-2pt]\footnotesize 平方项四倍};
\end{tikzpicture}
\end{center}

\emph{两条增长曲线对应两个场景。训练撞的是注意力矩阵的平方增长，在相对较短的 \(n\) 上就顶到显存天花板；推理撞的是 KV cache 的线性增长，斜率小得多，但它必须在整个生成过程中一直持有，并且按并发请求数成倍叠加。优化其中一条，另一条不会自动跟着好转。}

图里那条线性曲线的斜率看着温和，可它是每条请求各一份。服务端同时开 32 路对话，等效斜率就是图上的 32 倍——推理的墙来得比这张图看上去要快。

\subsection{四条缓解路线，各自卖掉了什么}

削减 \(n^2\) 的办法不少，每一条都在某个地方付了钱。把代价说清楚比记住名字有用。

\emph{稀疏与局部注意力}把每个查询的可读集合从全部 \(n\) 个位置缩到一个窗口或一个固定模式，计算量降到 \(O(nw)\)。卖掉的是任意两个位置单层直达的能力。相隔 500 步而窗口只有 256 时，信息必须靠中间层逐跳传递，而每一跳都要经过 softmax 加权和残差混合，长程依赖被稀释。这恰好抵消了下一篇要讲的 Transformer 最核心的结构优势，所以窗口大小不是越省越好，得对着任务真实的依赖距离定。

\emph{线性注意力}换了个算式的结合顺序。把 \(\exp(q^{\trans}k)\) 换成一个可分解的相似度 \(\phi(q)^{\trans}\phi(k)\)，注意力就写成

\[
\bigl(\phi(Q)\phi(K)^{\trans}\bigr)V \;=\; \phi(Q)\bigl(\phi(K)^{\trans}V\bigr).
\]

左边先算 \(n\times n\)，右边先算 \(\phi(K)^{\trans}V\)，这是个 \(d'\times d\) 的矩阵，与 \(n\) 无关。同一个结果，两种括号，复杂度从 \(O(n^2d)\) 掉到 \(O(nd d')\)。卖掉的是尖锐性。softmax 的指数让最大的那个分数指数级压过其余，注意力可以逼近一次精确检索——``把第 137 个位置的内容原样取过来''这类操作靠的就是这种尖锐。而 \(\phi\) 的输出维度 \(d'\) 有限，得到的相似度矩阵秩至多 \(d'\)，无法在任意位置制造出一根足够尖的峰，注意力被迫变平滑。需要精确复制和检索的任务上，线性注意力的退化最明显。这与 \hyperref[sec:03-Linear-Algebra/07-SVD-and-Data-Applications/03]{``低秩近似：用少数方向保留主要结构''}是同一件事的两面：低秩能保住主要结构，保不住尖峰。

\emph{MQA 与 GQA}只动缓存。多个 Query 头共享同一组 Key 和 Value——共享到底就是 MQA，分成 \(g\) 组各共享一份就是 GQA。缓存量除以 \(h\) 或 \(h/g\)，上面那个 16 GB 在 \(h=32\) 全共享时变成 500 MB。卖掉的是 \hyperref[sec:14-Deep-Learning/05-Attention-and-Sequence-Models/02]{``多头不是冗余而是分工''}里那套分工：Value 只剩一份，各头再怎么算权重，读回来的内容都落在同一个子空间里，能并存的对齐关系随之减少。GQA 取 \(g=8\) 一类的中间值，是因为经验上多数头的 Key/Value 确实高度相似，砍到一份太狠、保留 32 份太浪费——同样是经验规律。

\emph{FlashAttention 不在这份名单里}，尽管它经常被和上面三条并列。它一个数学符号都没改：算出来的注意力与朴素实现在浮点误差内相同，没有近似、没有稀疏、没有共享。它做的是分块。把 \(Q,K,V\) 切成能装进片上 SRAM 的小块，在块内完成分数、softmax 与加权求和，用一个在线更新最大值和归一化因子的技巧保证分块 softmax 与整行 softmax 结果一致，全程不把 \(n\times n\) 的中间矩阵写回显存。

区分两种复杂度，这一段才算读懂。\emph{算法复杂度}是浮点运算次数，FlashAttention 仍是 \(O(n^2d)\)，一次没省。\emph{访存复杂度}是显存与片上存储之间搬运的数据量，朴素实现要往返搬 \(O(n^2)\) 个数，FlashAttention 降到 \(O(n^2 d/M)\)（\(M\) 是 SRAM 容量）。上一小节说过注意力是访存受限的，所以砍访存直接换来数倍加速，同时显存占用从 \(O(n^2)\) 降到 \(O(n)\)——不再存那个大矩阵，反向传播时按需重算。

这个区分不是术语洁癖。它决定了 FlashAttention 能做什么、不能做什么：它把显存墙推远了很多，让 32k 从不可能变成可能；但计算量仍随 \(n^2\) 增长，\(n=10^6\) 时该慢还是慢。想突破平方，必须像前三条那样真的改变数学，也就必须真的卖掉点什么。

前四篇都在同一个框架内做加减法：注意力换来了内容决定的读取，为此欠下的是位置、深度稳定性和这一章的复杂度账。下一篇退后一步，把注意力和它取代的那条技术路线放在同一张表上比，看看放弃的东西是不是也该算进来。
